\documentclass[12pt]{article}
%\usepackage[OT1]{fontenc}
%\usepackage{babel}
%\usepackage{hyperref} 
%\usepackage{url}   %this allows us to cite URLs in the text
\usepackage{graphicx}  %allows for graphic to float when doing jou or doc style
\usepackage{amssymb}  %use formatting tools  for math symbols
\usepackage{natbib}  %might conflict with apa...
\usepackage{booktabs}
%\usepackage{Sweave}
%\usepackage{tikz}
%\usepackage{pgf}
\usepackage{fullpage}
\usepackage{setspace} 
%\usepackage{lineno}
%\linenumbers


\begin{document}  
\title{Statement of Contributions to Diversity}
\author{Solomon Messing}
\date{September 21, 2014}
\maketitle 

I try to contribute to diversity on campus in my research and teaching.  I have a number of published papers and works in-progress documenting subtle racial biases in modern political communication, in the hopes that with purposeful consideration of these biases, individuals can counteract them \citep[scientific evidence of this possibility can be be found in][]{devine_regulation_2002}.  In ``Social News in Black and White: How Ethnic and Socioeconomic Attributes Shape the Way We Consume and Evaluate Social Media,'' I examine how one highly visible attribute---group membership---affects how people select and evaluate media content in social media, where the attributes of the recommenders are increasingly salient.  Using an innovative survey experiment conducted with a national sample, I show that ethnic and socio-economic attributes affect decisions to read content, and that reading identical content---a story about welfare in this case---shared by different group members affects subsequently measured political preferences.  In particular, Whites come away with positive evaluations of welfare policies after reading an article recommended by counter-stereotypical Blacks, but tend to provide on balance negative evaluations when recommenders are Whites or stereotype-consistent Blacks.  

In ``Bias in the Flesh: Skin Complexion and Stereotype Consistency in Political Campaigns'' (with Ethan Plaut and Maria Jabon, \emph{forthcoming at Public Opinion Quarterly}) we extend findings linking skin complexion to negative stereotypes and adverse real-world outcomes to the realm of political ad campaigns, in which skin complexion can be easily manipulated in ways that are difficult to detect.  We devise a method to measure how dark a candidate appears in an image, and use it to examine how complexion varies with ad content during the 2008 campaign, finding that darker images were more frequent in negative ads---especially in those linking Obama to crime---which aired more frequently as election day approached.  We then conduct an experiment to document how these darker images can activate stereotypes, which shows that a subtle darkness manipulation is sufficient to activate the most negative stereotypes about Blacks---even when the candidate is a famous counter-stereotypical exemplar---Barack Obama.  We then collect further evidence of an evaluative penalty for darker skin based on an observational study measuring affective responses to depictions of Obama with varying skin complexion, presented via the Affect Misattribution Procedure in the 2008 American National Election Survey.  The work demonstrates that darker images are used in a way that compliments ad content, and shows that doing so can negatively affect how individuals evaluate candidates and think about politics.

The next step in this research agenda is to examine the relationship between skin complexion and electoral outcomes.  Toward that end, I've assembled an images database of each primary and general election candidate for national and some state office in 2010 and am merging this with other candidate data.  Preliminary evidence suggests a skin complexion penalty at the ballot box---the measure of skin complexion referenced above is strongly correlated with electoral victory among 2010 minority candidates for national and state-level office.  Future work will expand this database to other election years and examine implications for representation and policy making.  Furthermore, in experiments conducted with Shanto Iyengar, we show that presenting participants with darker images of a minority candidate decreases feelings of warmth and the probability of voting for the candidate in online experiments, especially for individuals who test high for negative implicit associations with Blacks.

In the course of living in South Africa, Indonesia, and Japan, and working closely with Chinese nationals on a daily basis, I have cultivated communication skills and cross-cultural experiences that have prepared me to maximize effective collaboration and teaching with a diverse community.  The most import facet of these skills is learning the right way to listen.  For me, this goes deeper than the helpful step of learning foreign languages, but includes trying to understand the cultural context---the unspoken rules, incentives, and assumptions that individuals grow up with and know.  Learning about this context can help unlock the pathway to motivate and help elucidate key skills and knowledge.

I have put these skills to use in the classroom at Stanford, teaching a class that examines the role of social technologies in modern politics and electoral campaigns.  The course included students from diverse backgrounds including student athletes and high school students, many from underprivileged groups and backgrounds.  Knowing that students from these groups only rarely take courses with statistics and programming content motivated me to ensure that the course was accessible to students from a variety of different backgrounds.  Though students often found this material difficult, I worked with students in groups and individually, sharing with them why this material is so exciting and highlighting the most important parts of it, they ultimately developed a solid introductory understanding of the material and reported a positive experience in the classroom. 

I am dedicated to contributing to the diversity of the campus community, especially in teaching and in research.  


\bibliographystyle{POQ}
\bibliography{Race}

\end{document}









